\begin{afterword}
\summaryhead

A short story, adapted from Raymond Smullyan. A court jester gives the king two boxes, one holding gold and the other an angry frog. Each box bears an inscription about the boxes and about which inscription is false, and the jester adds that exactly one inscription is true. The king opens the wrong box and is bitten. When the jester starts to explain the solution, the king has him thrown in the dungeon.

The next day the king shows the jester two boxes, one holding the key to his chains and the other a dagger. The first reads ``Either both inscriptions are true, or both inscriptions are false''; the second reads ``This box contains the key.'' The jester reasons that the second inscription is true whether the first is true or false, opens the second box, and finds the dagger. He cries that this is logically impossible. The king answers: ``I merely wrote those inscriptions on two boxes, and then I put the dagger in the second one.''

\responsehead

The story is fiction, and it is judged here for what it shows. What it shows is that reasoning about what is written on a box tells you about the box only if someone has promised that the writing is true. It shows this faithfully, and its two halves make the contrast. In the first puzzle the jester gives such a promise, and the puzzle has exactly one answer, which the king misses. In the second the king gives none, and the jester's valid reasoning fails.

The credit to Smullyan is accurate and, if anything, modest. The second puzzle has the structure of problem 70 in \textsc{What Is the Name of This Book?}, and the king's last line is close to Smullyan's own solution: anyone can put an object in a casket and write ``any inscriptions at all on the lids.'' The first half, which the post adds, dramatizes the contrast Smullyan draws with his earlier puzzles, where a promise about the inscriptions would have made a misplaced object mean that someone had lied.

What the story leaves out is the explanation. Smullyan gives two: the suitor had no information about the truth of the sentences, and he assumed that each sentence was either true or false. The second is the more instructive here. With the dagger in the second box, the first inscription can be neither true nor false, so the jester's cry, ``It's logically impossible!'', is right about his assumption and wrong about the boxes. The story gives only the king's line, which, read alone, could be taken as a verdict against logic. The next post, ``The Parable of Hemlock,'' makes a related point in its closing verse: ``Logic stays true, wherever you may go, / So logic never tells you where you live.''

In short: a faithful and credited adaptation of Smullyan's puzzle, whose two halves show that inscriptions tell you about the boxes only if someone has promised they are true, told without the explanation its source supplies.
\end{afterword}
